1-------------------
\subsection{Contexte et motivations}
% essor des produits dérivés, attrait du modèle binomial, lien actuariat / quante
\subsection{Terminologie}
% actif financiers, produits dérivés, valorisation, taux sans risque
\subsection{Arbitrage}
% Absence d'Opportunité d'Arbitrage intuitive
\subsection{Enoncé du problème}
% complexité du temps continu, problème de la discrétisation
\subsection{Objectifs}
%analyse théorique, développement informatique, validation et extension, analyse comparative

transition vers le modèle de Black-Scholes

4-------------------------------------------
Conception et implémentation en C$\texttt{++}$ du modèle CRR
    \item Architecture (encapsulation/gestion de la mémoire)
    \item algorithme principal
    \item complexité
    \item tests


5---------------------------------
\section{Convergence et valorisation d'options complexes}
\begin{itemize}
    \item Portefeuille continu
    \item options américaines
    \item options à barriere
    \item calcul des grecques
\end{itemize}


6-----------------------------
\section{Comparaisons avec d'autres modèles}
\begin{itemize}
    \item comparaison avec black scholes
    \item comparaison avec monte carlo
    \item comparaison avec un arbre trinomial
\end{itemize}



7------------------------
\section{Analyse des limites et possibilités d'extension}
\begin{itemize}
    \item cas réel
    \item limites
    \item Améliorations possibles : arbres recombinants, modèles multi-actifs, options exotiques
\end{itemize}